% Part: first-order-logic
% Chapter: semantics
% Section: intro-semantics

\documentclass[../../../include/open-logic-section]{subfiles}

\begin{document}

\olfileid{fol}{syn}{its}

\olsection{پېژندنه}

عبارتونو ته معنا ورکول د معنٰی پوهنې کار دے۔ د معنٰی پوهنې مرکزي
مفهوم په !!a{structure} کښې د صدق مفهوم دے۔ !!^a{structure} د ژبې
بنسټيزو اجزاوو ته معنا ورکوي: !!a{domain} د څيزونو يو ناتش سټ دے۔
سورونه داسې تفسيرېږي چې پر همدې ساحه ګرځي، !!{constant}s د ساحې
له غړو سره ګومارل کېږي، !!{function}s له !!{domain} څخه
همدې ساحې ته له تابعو سره ګومارل کېږي، او !!{predicate}s پر
!!{domain} له اړيکو سره ګومارل کېږي۔ !!{domain} او د بنسټيز لغت
ګومارنې په ګډه !!a{structure} جوړوي۔ !!^{variable}s په !!{formula}s
کښې راتلے شي، او د معنٰی پوهنې د ورکولو دپاره بايد هغو ته هم د
!!{domain} !!{element}s وګومارو---دا د متغير ګومارنه ده۔ په پای کښې
د صدق اړيکه دا ټول سره يوځاے کوي۔ !!^a{formula} په
!!a{structure}~$\Struct{M}$ کښې د !!a{variable} ګومارنې~$s$ په نسبت
صدق کولے شي، چې $\Sat{M}{!A}[s]$ ليکل کېږي۔ دا اړيکه هم د~$!A$ پر
جوړښت په استقرا تعريفېږي: د منطقي نښلوونکو د رښتياوالي جدولونه
کاروو، څو مثلاً د $(!A \land !B)$ صدق د $!A$ او~$!B$ د صدق يا نه
صدق له مخې تعريف کړو۔ بيا څرګندېږي چې که !!{formula}~$!A$
!!a{sentence} وي، يعنې ازاد متغيرونه ونۀ لري، د !!{variable} ګومارنه
اړونده نۀ ده؛ نو ويلے شو چې !!{sentence}s په ساده ډول په
!!{structure}s کښې صدق کوي يا نۀ کوي۔

د !!{sentence}s دپاره د صدق د اړيکې $\Sat{M}{!A}$ پر بنسټ بيا د
اعتبار، معنايي استلزام، او د صدق وړتيا بنسټيز معنايي مفهومونه
تعريفولے شو۔ !!^a{sentence} هله معتبره ده، $\Entails !A$، چې هر
!!{structure} پکښې صدق وکړي۔ د !!{sentence}s يو سټ هله هغه لازموي،
$\Gamma \Entails !A$، چې هر هغه !!{structure} چې د~$\Gamma$ په ټولو
!!{sentence}s کښې صدق کوي، په~$!A$ کښې هم صدق وکړي۔ او د
!!{sentence}s يو سټ هله د صدق وړ دے چې کوم !!{structure} په هغه کښې
په يو وخت په ټولو !!{sentence}s کښې صدق وکړي۔ څرنګه چې !!{formula}s
په استقرايي ډول تعريف شوي، او صدق بيا د !!{formula}s پر جوړښت په
استقرا تعريفېږي، استقرا کارولے شو څو د خپلې معنٰی پوهنې خاصيتونه ثابت
او تعريف شوي معنايي مفهومونه سره وتړو۔

\end{document}
